\chapter{Software Requirement Specifications}
\label{ch:srs}

This chapter translates the reviewed problem and evidence into a specification for the Manzil research prototype. The requirements describe observable behavior rather than implying that every item has already been implemented. Where the available repository evidence does not support a measured threshold, the text states a design target or requires profiling before a value is fixed. The specification gives particular attention to reproducibility, legal signal transitions, observation freshness, route-version consistency, and visible failure handling because these properties determine whether later experiments can be interpreted. It also separates the video-perception workflow from SUMO-based control experiments so that estimated observations are not silently presented as simulator ground truth. Account management, public dispatch integration, traffic enforcement, clinical response prediction, and physical signal deployment remain outside the approved scope. The resulting requirements provide a contract for the architecture in Chapter \ref{ch:design} and a traceable basis for the evaluation procedures in Chapter \ref{ch:method}.

\section{List of Features}

The prototype is organized into six feature groups that follow the operator's experimental workflow from configuration to evidence export. The grouping prevents a user-interface screen or implementation library from being mistaken for a requirement and clarifies which service owns each outcome. It also shows that emergency operation is not an isolated routing feature: a valid mission depends on scenario state, traffic observations, signal control, monitoring, and durable records. The following statements summarize those groups before the individual functional requirements define testable behavior.

\begin{enumerate}
\item The operator can select a network, demand profile, controller, seed, incident, duration, and optional emergency mission before a run begins.
\item The perception workflow can process licensed recorded video into detections, tracks, road-region observations, and timestamped aggregate measures with confidence information.
\item The simulation workflow advances SUMO through one owner, reads traffic state, applies fixed or adaptive requests, and preserves legal phase transitions.
\item The emergency workflow validates endpoints, computes a dynamic path, applies stable rerouting, activates bounded moving priority, detects passage, and recovers normal control.
\item The monitoring interface presents map state, signals, traffic indicators, ambulance route, corridor state, warnings, and the current data-source mode.
\item The experiment service retains immutable configuration, events, metrics, software versions, seeds, and termination state for comparison and export.
\end{enumerate}

\section{Functional Requirements}

Table \ref{tab:functional-requirements} defines the mandatory system behavior using stable identifiers that are reused in the traceability matrix. The requirements distinguish a requested signal action from the action actually applied, require route and observation versions to be retained, and make fallback behavior part of the specification rather than an informal implementation detail. Several entries use the word ``shall'' because they are acceptance conditions for the proposed prototype; they are not claims of completed implementation. Together, FR-01 through FR-08 cover experiment configuration and normal control, FR-09 through FR-14 define the emergency-route and corridor lifecycle, and FR-15 through FR-20 cover monitoring, evidence, comparison, export, and controller failure. This arrangement lets tests target a coherent workflow while preserving individual identifiers for defect tracking. A requirement is satisfied only when its normal and failure paths are demonstrated with retained evidence. For example, returning a path is insufficient for FR-10 if the cost snapshot and graph-validity conditions are not also recorded.

\begin{longtable}{p{0.10\textwidth}p{0.78\textwidth}}
\caption{Functional requirements}\label{tab:functional-requirements}\\
\toprule
\textbf{ID} & \textbf{Requirement}\\
\midrule\endfirsthead
\toprule
\textbf{ID} & \textbf{Requirement}\\
\midrule\endhead
FR-01 & The system shall validate a scenario configuration before starting and shall reject missing network, controller, seed, duration, or demand inputs.\\
FR-02 & The system shall start, pause, resume, and stop one controlled SUMO run while preserving an unambiguous run state.\\
FR-03 & The simulation adapter shall be the sole owner of SUMO stepping and signal actuation.\\
FR-04 & The perception worker shall process a selected recorded video and emit timestamped detections, tracks, source metadata, and processing status.\\
FR-05 & The observation service shall aggregate valid vehicle observations into defined road, lane, or approach intervals without silently mixing sources.\\
FR-06 & The operator shall be able to select fixed-time, rule-adaptive, or the configured RL controller before a run.\\
FR-07 & Every signal action shall identify its intersection, requested and applied phase, effective simulation step, controller, and reason.\\
FR-08 & The system shall enforce minimum green and configured yellow/all-red clearance before incompatible movements receive right of way.\\
FR-09 & The emergency service shall validate ambulance origin and destination and return a clear no-route response if either is invalid or disconnected.\\
FR-10 & The route service shall compute a valid remaining path using nonnegative time-based edge costs and shall record the cost snapshot.\\
FR-11 & The rerouting service shall re-evaluate the remaining path after a closure, invalidated edge, or sustained cost improvement above the configured threshold.\\
FR-12 & The corridor coordinator shall activate priority only for a bounded look-ahead set of intersections on the current ambulance route.\\
FR-13 & The corridor coordinator shall detect passage or cancellation and restore the affected intersection to a safe normal-control state.\\
FR-14 & A route revision shall cancel priority requests that no longer belong to the remaining path before activating new ones.\\
FR-15 & The dashboard shall display the current source mode, simulation status, controller, signal state, traffic indicators, ambulance position, route, and corridor state.\\
FR-16 & The dashboard shall present stale, invalid, disconnected, and fallback states visibly rather than continuing with apparently current data.\\
FR-17 & The system shall record configuration, seed, component versions, periodic metrics, route revisions, signal events, failures, and termination reason for each run.\\
FR-18 & The comparison service shall align selected runs by metric definition and warn when network, demand, duration, warm-up, or measurement versions are incompatible.\\
FR-19 & The operator shall be able to export run configuration and metrics in a documented machine-readable format.\\
FR-20 & On loss of the learned controller or invalid action, the control service shall fall back to a configured safe baseline and record the transition.\\
\bottomrule
\end{longtable}

\section{Quality Attributes}

The quality model for Manzil follows the sample report's separation between broad software qualities and individually testable non-functional requirements. The qualities below explain why the requirements matter to this research prototype, while Section 4.4 assigns stable identifiers to the corresponding acceptance conditions. The distinction is important because a feature may execute without producing credible research evidence. For example, a controller can return an action and still violate reliability if the applied phase, effective step, or fallback reason is lost. Similarly, a dashboard can appear responsive while showing stale observations. Manzil therefore treats usability, reliability, security, performance, maintainability, extensibility, and traceability as constraints on the complete workflow rather than as cosmetic objectives. Each quality is interpreted within a simulation and recorded-video setting; none is presented as evidence of certified public-road operation.

\subsection{Usability}

The interface is designed for a trained research operator who must configure a scenario, observe a run, initiate a simulated ambulance mission, and inspect retained evidence without switching among unrelated tools. Inputs use project terminology and display units, validation state, data source, and freshness beside the value that they qualify. Navigation preserves the selected run and mission context so that moving between the network overview, perception detail, emergency route, configuration, and analytics screens does not silently change the experiment. Status is not communicated by color alone: labels and symbols distinguish normal, stale, fallback, invalid, and disconnected states. Destructive or irreversible actions, such as stopping a run or replacing an active route version, require clear confirmation or are mediated by backend validation. These provisions do not claim that a formal usability study has been completed. They define observable interface behavior that can later be checked through task-based review using the proposed screens in Section \ref{sec:gui}.

\subsection{Reliability and Integrity}

Reliability concerns the preservation of valid state and evidence when a component is slow, unavailable, or produces an invalid request. One orchestrator and one SUMO adapter own simulation stepping and signal actuation, which prevents two services from advancing the same clock or issuing competing phase changes. Requested actions are stored separately from applied actions so the safety filter's intervention remains visible. A run configuration becomes immutable after start, route revisions are monotonically versioned, and corridor requests retain the route version that authorized them. Missing observations remain missing rather than being converted to zero, while stale observations are rejected or handled through a declared fallback. Terminal run states include completed, stopped, failed, and incomplete, each with a reason. Integrity therefore means that an unsuccessful experiment remains interpretable; it does not mean that every component is continuously available. Recovery and fallback behavior are specified explicitly in the use cases, risk register, and Chapter \ref{ch:design}.

\subsection{Security and Confidentiality}

Manzil is not a public dispatch or traffic-control service, but its development data still require controlled handling. Recorded video must be licensed, public under suitable terms, or specifically authorized for team use. Raw video and model artifacts remain external files referenced through provenance and checksums rather than being copied indiscriminately into experiment exports. Machine credentials, private paths, and access tokens are excluded from run manifests. The proposed dashboard is an operator interface for a controlled laboratory environment and does not introduce public user registration merely because sample reports contain account modules. If authentication is later required for deployment, it must be added as a separate verified requirement rather than assumed in this deliverable. Exported research records should contain scenario and measurement identifiers but no unnecessary personal data. These boundaries provide proportionate confidentiality for the stated prototype while avoiding unsupported claims about production authorization, encryption infrastructure, or compliance with operational emergency-service policies.

\subsection{Performance}

Performance is defined relative to the chosen simulation step, control interval, video workload, and demonstration hardware. The control path must complete its decision before the configured deadline for the step in which the action becomes effective; otherwise the interval is marked late and a declared baseline or safe continuation is used. Video inference runs in a bounded worker so a slow model cannot block simulation stepping or an API request indefinitely. Dashboard delivery is a lossy view of authoritative state, allowing an inactive client to skip intermediate frames without delaying control. Route computation and comparison queries are profiled separately from simulation actuation because their acceptable latency depends on network size and update frequency. The repository contains no measured benchmark for the final hardware, detector, or road network, so exact thresholds remain design targets to be set after baseline profiling. This treatment prevents aspirational response times from being presented as achieved results and ensures that later measurements identify workload, hardware, and software versions.

\subsection{Maintainability and Reusability}

Maintainability is supported through typed, versioned contracts and replaceable implementations behind stable responsibilities. Fixed-time, rule-adaptive, and learned controllers share one request interface and therefore receive the same snapshot fields and legal phase constraints. The repository boundary allows early file-backed persistence to be replaced by PostgreSQL without changing the meaning of a run, while the perception boundary permits a detector or tracker to change only when its model version and output semantics are declared. Configuration is stored in validated files rather than scattered constants, and public interfaces are documented with units, identifiers, timing semantics, and failure behavior. Tests target cost functions, route continuity, freshness rules, schema validation, and corridor invariants independently of the final dashboard. Reuse is limited to components whose assumptions remain valid; a model, road network, or phase plan is not treated as reusable merely because its file can be loaded. These rules reduce coupling while preserving the provenance required for research comparison.

\subsection{Extensibility and Traceability}

The design leaves controlled extension points for additional controllers, perception models, route-cost providers, scenarios, and visualizations. An extension must implement the existing contract or introduce an explicit schema version, and it must not acquire direct TraCI access or bypass the common safety filter. New measurements include a definition, unit, interval, and source run identifiers so previous and revised metrics cannot be combined silently. Traceability connects functional requirements to use cases, screens, architecture components, events, and planned tests through stable identifiers and cross-references, even though the standalone traceability matrix was removed at the supervisor's direction. Every exported comparison retains the run manifest, seed coverage, controller version, metric version, and failed-run status. This approach allows the prototype to grow without changing the interpretation of earlier evidence. Extensibility is therefore not an unrestricted plugin claim; it is the ability to replace or add a component while declaring the effect on control ownership, data contracts, and reproducibility.

\section{Non-Functional Requirements}

Functional behavior alone would permit a demonstration that appears successful but cannot be reproduced or audited. Table \ref{tab:nfr} therefore states the quality attributes that constrain every feature group. Reproducibility and data integrity require a run to preserve the exact scenario, version, seed, and metric definition used to generate a result. Safety requirements apply within simulation and prevent controllers from bypassing clearance or ownership rules. Freshness and isolation ensure that a stale perception result or slow dashboard cannot silently influence a later simulation step. Responsiveness is intentionally expressed as a profiled design target because the repository contains no defensible hardware benchmark from which to invent a deadline. Privacy is limited to the handling of licensed or team-authorized video and does not imply a production governance framework. These non-functional requirements also define how negative outcomes are handled: an unavailable controller, disconnected simulator, or incomplete record must remain visible rather than being removed from the evaluation set without explanation.

\begin{longtable}{p{0.11\textwidth}p{0.22\textwidth}p{0.57\textwidth}}
\caption{Non-functional requirements and design targets}\label{tab:nfr}\\
\toprule
\textbf{ID} & \textbf{Attribute} & \textbf{Requirement or target}\\
\midrule\endfirsthead
\toprule
\textbf{ID} & \textbf{Attribute} & \textbf{Requirement or target}\\
\midrule\endhead
NFR-01 & Reproducibility & A run shall store network and demand identifiers, seed, controller configuration, software/model versions, warm-up, duration, and metric schema.\\
NFR-02 & Safety in simulation & Incompatible green movements shall never be applied; every ordinary and emergency transition shall pass the safety filter and use the configured clearance sequence.\\
NFR-03 & Auditability & Applied decisions shall be distinguishable from requests, and their simulation timestamps and reasons shall be retained.\\
NFR-04 & Freshness & Each observation shall carry produced-at and effective simulation time. Stale thresholds are configuration values and stale data shall not silently drive control.\\
NFR-05 & Responsiveness & The control action for step $t$ should be available before its configured decision deadline. The exact deadline is a design target to be set after baseline profiling, not an achieved value.\\
NFR-06 & Isolation & Background video inference and dashboard clients shall not advance SUMO or directly actuate signals.\\
NFR-07 & Usability & A trained operator should be able to configure the core demonstration from one experiment screen, with units and validation messages visible.\\
NFR-08 & Portability & Services should be runnable from versioned environments; containerization is proposed after local baseline integration is stable.\\
NFR-09 & Maintainability & Component interfaces shall use typed, versioned data contracts and controller implementations shall share one interface.\\
NFR-10 & Data integrity & A completed run configuration shall be immutable; derived summaries shall retain their metric version and source run identifiers.\\
NFR-11 & Privacy & Video shall be limited to licensed or team-authorized material, stored only as needed, and excluded from exports unless explicitly requested.\\
NFR-12 & Availability & The prototype shall fail visibly and recover to a safe stopped or baseline-controlled simulation state; no public-service availability is claimed.\\
\bottomrule
\end{longtable}

\section{Assumptions}

SUMO and TraCI are required for controlled traffic experiments \cite{lopez2018sumo,sumodocs}. A selected network must contain legal lane connections and signal plans after import and manual cleanup, and representative routes must be validated before it is used for comparison. Video sources require sufficient visibility, a documented license or team authorization, and annotations appropriate to the selected measure. The simulation assumes one clock owner, nonnegative route costs, and emergency endpoints on ambulance-accessible edges. It also assumes that configured phase transitions reflect the validated simulation network; they are not certified timings for physical signal hardware. The dashboard is an operator and evaluator tool rather than an autonomous public dispatch system, and a displayed route does not constitute operational navigation advice. When an assumption cannot be met, the affected workflow must stop or disclose a bounded fallback. It must not silently substitute a different network, data source, time unit, or observation mode because that would make later comparisons irreproducible.

React and FastAPI are fixed by the approved proposal for the interface and application service boundary \cite{reactdocs,fastapidocs}. Python is implied by the planned perception, reinforcement-learning, and TraCI ecosystem. PostgreSQL, a WebSocket transport, Docker, and specific chart or map libraries are proposed dependencies and remain replaceable until implementation confirms them. These assumptions are deliberately narrower than a production deployment specification. Manzil is evaluated as a controlled prototype on prepared networks, not as an interface to public signal controllers or ambulance dispatch infrastructure. A selected video source must have suitable permission and sufficient visibility for the intended annotation task, while the simulation network must be checked for lane connectivity, valid traffic-light programs, and ambulance-accessible routes. If any assumption fails, the run must be rejected or marked with an explicit limitation instead of substituting an undocumented configuration. This treatment keeps dependencies traceable and prevents a convenient library choice from becoming an unsupported architectural claim.

\section{Use Cases}

The primary actor is a research operator who configures experiments, observes their progress, initiates a simulated emergency mission, and exports evidence. SUMO, recorded video, and the experiment store support those actions but do not independently choose research parameters. Evaluators consume dashboard views and exported records without receiving control privileges. Figure \ref{fig:use-case} presents this boundary and connects the operator to eight goal-oriented interactions. Configuration and simulation appear separately because validation must finish before the single clock owner starts a run. Dispatch, corridor monitoring, and rerouting are also distinct so that a route revision cannot be treated as an invisible side effect. Tables \ref{tab:uc1}--\ref{tab:uc8} then specify a summary, conditions, special requirements, the basic actor--system exchange, and alternative behavior for each interaction. Collectively, they test more than the successful interface path: invalid endpoints, stale streams, no-route cases, controller failure, incompatible comparisons, and termination reasons must all produce observable outcomes. The use cases therefore define acceptance evidence as well as operator interaction.

\begin{figure}[htbp]
\centering
\begin{tikzpicture}[font=\footnotesize,uc/.style={draw,ellipse,minimum width=3.5cm,minimum height=0.8cm,align=center},node distance=5mm and 11mm]
\coordinate (op) at (0,0);
\draw (op) ++(0,7mm) circle (2mm);
\draw (op) ++(0,5mm)--++(0,-7mm);
\draw (op) ++(-4mm,2mm)--++(8mm,0);
\draw (op) ++(0,-2mm)--++(-4mm,-5mm);
\draw (op) ++(0,-2mm)--++(4mm,-5mm);
\node[below=8mm of op,align=center] {Research\\operator};
\node[uc,right=20mm of op,yshift=18mm] (cfg) {Configure scenario};
\node[uc,below=of cfg] (run) {Run/control simulation};
\node[uc,below=of run] (video) {Process traffic video};
\node[uc,below=of video] (inspect) {Inspect traffic state};
\node[uc,right=17mm of cfg] (dispatch) {Dispatch ambulance};
\node[uc,below=of dispatch] (corr) {Monitor route/corridor};
\node[uc,below=of corr] (reroute) {Handle route change};
\node[uc,below=of reroute] (compare) {Compare/export runs};
\node[draw,fit=(cfg)(run)(video)(inspect)(dispatch)(corr)(reroute)(compare),inner sep=5mm,label={[anchor=north west]north west:Manzil research prototype}] (boundary) {};
\foreach \x in {cfg,run,video,inspect}{\draw (op)--(\x);}
\draw (op) -- ++(5mm,0) |- (dispatch.west);
\draw (op) -- ++(5mm,0) |- (corr.west);
\draw (op) -- ++(5mm,0) |- (reroute.west);
\draw (op) -- ++(5mm,0) |- (compare.west);
\end{tikzpicture}
\caption{Proposed operator use-case diagram}
\label{fig:use-case}
\end{figure}

\newcommand{\uctable}[8]{%
\begin{table}[htbp]\centering
\caption{#2}\label{#1}
\renewcommand{\arraystretch}{1.05}
\begin{tabularx}{\textwidth}{|p{0.06\textwidth}|p{0.34\textwidth}|p{0.06\textwidth}|Y|}
\hline
\multicolumn{2}{|l|}{Name} & \multicolumn{2}{l|}{#2}\\ \hline
\multicolumn{2}{|l|}{Actors} & \multicolumn{2}{l|}{Research operator}\\ \hline
\multicolumn{2}{|l|}{Summary} & \multicolumn{2}{p{0.50\textwidth}|}{#3}\\ \hline
\multicolumn{2}{|l|}{Pre-Conditions} & \multicolumn{2}{p{0.50\textwidth}|}{#4}\\ \hline
\multicolumn{2}{|l|}{Post-Conditions} & \multicolumn{2}{p{0.50\textwidth}|}{#5}\\ \hline
\multicolumn{2}{|l|}{Special Requirements} & \multicolumn{2}{p{0.50\textwidth}|}{#6}\\ \hline
\multicolumn{4}{|c|}{Basic Flow}\\ \hline
\multicolumn{2}{|c|}{Actor Action} & \multicolumn{2}{c|}{System Response}\\ \hline
#7
\multicolumn{4}{|c|}{Alternative Flow}\\ \hline
#8
\hline
\end{tabularx}\end{table}}

\uctable{tab:uc1}{UC-01: Configure a Scenario}{The operator prepares a validated and versioned experiment configuration.}{Network, demand, controller, and source descriptors are available.}{A valid configuration is stored; no run has started.}{Identifiers, durations, seeds, and controller compatibility must be validated.}{1 & Opens experiment configuration. & 2 & Displays available networks, demand profiles, controllers, and source modes.\\ \hline 3 & Selects settings and submits. & 4 & Validates the fields and stores the versioned configuration.\\ \hline}{3 & Submits missing or incompatible values. & 4-A & Identifies each invalid field and does not create a runnable configuration.\\}

\uctable{tab:uc2}{UC-02: Run and Control a Simulation}{The operator starts and monitors a configured SUMO experiment.}{UC-01 is valid and the SUMO adapter is available.}{The run is completed, stopped, or failed with an explicit terminal reason.}{Only the orchestrator may advance simulation time; illegal actions require fallback.}{1 & Selects Start for a valid configuration. & 2 & Creates the run, launches SUMO, and completes warm-up.\\ \hline 3 & Monitors, pauses, resumes, or stops the run. & 4 & Steps the simulation, applies lawful actions, records evidence, and terminates safely.\\ \hline}{1 & Starts an unavailable or invalid setup. & 2-A & Rejects the start or records launch/TraCI failure while retaining diagnostic logs.\\}

\uctable{tab:uc3}{UC-03: Process Traffic Video}{The operator creates a video-processing job that produces traceable traffic observations.}{A licensed video source and model configuration are available.}{Stored observations identify the source, model, interval, units, confidence, and job status.}{Uncalibrated video may report proxies but not physical queue length or speed.}{1 & Selects a video source and model configuration. & 2 & Validates provenance, format, model, and optional calibration.\\ \hline 3 & Starts processing. & 4 & Detects, tracks, assigns regions, aggregates observations, and records job progress.\\ \hline}{3 & Submits unreadable media or processing fails. & 4-A & Reports the failure and preserves completed evidence without inventing physical measures.\\}

\uctable{tab:uc4}{UC-04: Inspect Traffic State}{The operator reviews current or recorded traffic state without changing control decisions.}{A simulation run or perception job has emitted observations.}{The displayed state remains traceable to its source and no decision is changed.}{Source mode, units, freshness, confidence, and invalid values must remain visible.}{1 & Opens the overview or an intersection detail. & 2 & Displays source, freshness, counts, queue values, phases, controller, and warnings.\\ \hline 3 & Selects another time or intersection. & 4 & Retrieves the corresponding traceable snapshot.\\ \hline}{1 & Opens a disconnected or stale stream. & 2-A & Marks the state stale or unavailable and does not present the last value as current.\\}

\uctable{tab:uc5}{UC-05: Dispatch an Ambulance}{The operator creates one simulated emergency mission between valid network positions.}{An emergency-enabled scenario is running and no core mission is active.}{A mission is active with a valid route, or the request is rejected cleanly.}{Origin and destination must be legal vehicle edges and corridor activation must be atomic.}{1 & Selects origin and destination and requests dispatch. & 2 & Validates the endpoints and computes a time-based legal route.\\ \hline 3 & Confirms the valid mission. & 4 & Spawns the ambulance, publishes the route version, and requests bounded look-ahead.\\ \hline}{1 & Submits invalid endpoints, no-route input, or a duplicate mission. & 2-A & Explains the rejection and activates no partial corridor.\\}

\uctable{tab:uc6}{UC-06: Monitor Route and Corridor}{The operator follows ambulance movement, route versions, and intersection priority states.}{An emergency mission is active.}{Mission progress and every priority transition are recorded.}{Priority must obey legal transition, expiry, passage, timeout, and recovery rules.}{1 & Opens the active emergency mission. & 2 & Displays position, remaining route, current route version, and corridor states.\\ \hline 3 & Continues monitoring as the ambulance advances. & 4 & Requests safe priority, detects passage, releases ownership, and records recovery.\\ \hline}{3 & A position becomes stale, a request is unsafe, or timeout occurs. & 4-A & Cancels affected requests, applies fallback or recovery, and displays an alert.\\}

\uctable{tab:uc7}{UC-07: Handle a Blockage or Route Change}{The system evaluates rerouting from the ambulance's valid current position.}{An active mission exists and a closure or sustained cost change is detected.}{Exactly one route version remains current and every decision retains its cause and cost snapshot.}{Feasibility, improvement threshold, persistence, cooldown, and obsolete-request cancellation apply.}{1 & Reviews the blockage or reroute notification. & 2 & Computes and validates a candidate route under one cost snapshot.\\ \hline 3 & Continues the mission after the decision. & 4 & Publishes an accepted version atomically or retains the current route with a reason.\\ \hline}{3 & No feasible route exists. & 4-A & Cancels future priority, marks the mission blocked, and records the graph snapshot.\\}

\uctable{tab:uc8}{UC-08: Compare and Export Runs}{The operator compares eligible experiments and exports their evidence and manifests.}{At least two completed or failed run records exist.}{The comparison and export identify source run IDs, metric definitions, and schema version.}{Network, demand, warm-up, duration, seed coverage, and metric versions must be disclosed.}{1 & Selects runs and a metric. & 2 & Checks comparability and displays distributions, variability, failures, and configuration differences.\\ \hline 3 & Requests an export. & 4 & Packages selected records with identifiers, units, definitions, and manifests.\\ \hline}{1 & Selects incompatible runs. & 2-A & Warns or refuses direct comparison and never silently omits failed runs.\\}

\section{Hardware and Software Requirements}

The development environment must support two different workloads: repeated model or controller experiments and a prepared, repeatable demonstration. The distinction avoids presenting a high-end training machine as a mandatory deployment requirement. Hardware values in this section are capability statements because the selected Lahore network, final detector, training budget, and demonstration machine have not yet been fixed. Software choices are separated from hardware so that a library can be versioned or replaced without changing the resource description. Tables \ref{tab:hardware} and \ref{tab:software} identify the responsibilities and confirmation points that must be resolved before implementation evidence is reported. The demonstration must use the same network, controller configuration, metric schema, and software versions as the retained evaluation record even if it runs fewer scenarios. Measured CPU, GPU, memory, storage, and response-time values will replace the current planning statements only after the representative workload has been profiled.

\subsection{Hardware Requirements}

Table \ref{tab:hardware} separates development and training resources from the minimum capabilities needed to reproduce a prepared demonstration. A multi-core processor is required for SUMO, service processes, preprocessing, and analysis. A CUDA-capable GPU is desirable for detector experiments and reinforcement-learning training, but its model and memory capacity remain unconfirmed; it is not required for every dashboard or SUMO-only workflow. Storage must accommodate approved video, virtual environments, checkpoints, network files, traces, and repeated run records while preserving checksums and versions. The demonstration system needs enough memory and compute to run the selected network, API, dashboard, and prepared policies without unbounded queues. If live video inference cannot meet the profiled budget, recorded inference output may be demonstrated as an explicitly labelled fallback. These conditions keep resource claims honest and allow the team to scale the network or model before committing to a numerical minimum.

\begin{table}[htbp]
\centering
\caption{Hardware resources for development and demonstration}
\label{tab:hardware}
\begin{tabularx}{\textwidth}{p{0.20\textwidth}YY}
\toprule
\textbf{Resource} & \textbf{Development and training} & \textbf{Prepared demonstration}\\
\midrule
Processor & Multi-core CPU for SUMO, preprocessing, services, and repeated analysis & CPU able to run the selected SUMO network, API, and dashboard at the configured step rate\\
Graphics processor & CUDA-capable GPU desirable for detector and RL experiments; exact model and memory are unconfirmed & Required only when live inference is included; otherwise prepared outputs or CPU-safe models may be used and labelled\\
Memory & Capacity for concurrent simulator, worker, services, model, and analysis processes; amount to be profiled & Capacity for the selected scenario and services without unbounded swapping or queue growth\\
Storage & Approved video, environments, checkpoints, maps, traces, exports, and backups & Versioned scenario, required models, application build, and demonstration records\\
Display/network & Development display and local networking for dashboard and service testing & Display suitable for the dashboard; local connectivity only unless a reviewed deployment requires more\\
\bottomrule
\end{tabularx}
\end{table}

\subsection{Software Requirements}

Table \ref{tab:software} follows the sample report's responsibility-oriented presentation: each technology is tied to a task and to a configuration item that must be recorded. SUMO and TraCI provide the authoritative simulation environment and control interface. Python supports perception, routing, control, experiment orchestration, and analysis, while React and FastAPI form the proposed interface boundary from the approved proposal. A detector, tracker, and reinforcement-learning library will be selected and versioned during implementation; the report does not claim a final checkpoint or tracker configuration. OpenStreetMap tools and SUMO NetEdit support network preparation, with a documented synthetic network retained as fallback. PostgreSQL remains proposed for structured persistence, but a file-backed repository may be used during integration if it preserves the same schema semantics. Git and environment lock files identify the software state used for every retained experiment.

\begin{table}[htbp]
\centering
\caption{Software technologies and configuration responsibilities}
\label{tab:software}
\begin{tabularx}{\textwidth}{p{0.24\textwidth}p{0.36\textwidth}Y}
\toprule
\textbf{Technology or resource} & \textbf{Responsibility} & \textbf{Configuration status}\\
\midrule
Python & Orchestration, perception, control, routing, metrics, and analysis & Interpreter and package lock recorded per experiment release\\
SUMO and TraCI & Microscopic traffic simulation, state access, route application, and signal actuation & Version and validated network package recorded\\
React & Operator dashboard and evidence views & Approved proposal choice; component and build versions to be recorded\\
FastAPI & Command, validation, query, and streaming boundary & Approved proposal choice; API schema versioned\\
Detector, tracker, and RL libraries & Vehicle observation and learned-controller experiments & Final frameworks, checkpoints, and hyperparameters remain to be confirmed\\
OpenStreetMap tools and SUMO NetEdit & Network extraction, cleanup, connection, route, and signal validation & Exact Lahore extract remains to be selected; synthetic fallback retained\\
Repository and persistence tools & Version control, manifests, structured records, and artifacts & Git required; PostgreSQL proposed behind a replaceable repository interface\\
\bottomrule
\end{tabularx}
\end{table}

\section{Graphical User Interface}
\label{sec:gui}

The interface specification follows the sample report by defining the screen inventory and navigation before presenting individual wireframes. The five views in Table \ref{tab:gui-screens} are proposed screens rather than screenshots of implemented software. They validate information hierarchy, source disclosure, error visibility, and the actions needed for a controlled demonstration without fabricating implementation progress. The configuration view establishes an immutable run, the overview supports monitoring, the perception view explains one intersection, the emergency view exposes route and corridor state, and the analytics view supports comparison after completion. Figure \ref{fig:gui-navigation} shows the permitted navigation among these views. Monitoring pages do not acquire control authority: every command still passes through the validated API boundary. This separation lets a reader understand both the interface flow and the evidence available at each stage while preserving the project status stated elsewhere in the report.

\begin{table}[htbp]
\centering
\caption{Proposed graphical user interface screen inventory}
\label{tab:gui-screens}
\begin{tabularx}{\textwidth}{p{0.10\textwidth}p{0.31\textwidth}p{0.20\textwidth}Y}
\toprule
\textbf{ID} & \textbf{Screen} & \textbf{Primary user} & \textbf{Purpose}\\
\midrule
G01 & Experiment Configuration & Operator & Validate and start a versioned scenario and controller configuration.\\
G02 & Traffic Overview & Operator and evaluator & Monitor network state, warnings, controller source, and active mission.\\
G03 & Perception and Intersection Detail & Operator and evaluator & Inspect observation provenance, freshness, confidence, and approach-level measures.\\
G04 & Emergency Dispatch and Route & Operator & Create a mission and monitor route versions and corridor transitions.\\
G05 & Comparative Analytics & Operator and evaluator & Check comparability, inspect metrics and failed runs, and export evidence.\\
\bottomrule
\end{tabularx}
\end{table}

\begin{figure}[htbp]
\centering
\begin{tikzpicture}[font=\footnotesize,node distance=9mm and 10mm,
screen/.style={draw,rounded corners,fill=manzilgray,minimum width=3.6cm,minimum height=.85cm,align=center},
arr/.style={-{Latex},thick}]
\node[screen] (config) {G01 Experiment\\Configuration};
\node[screen,below=of config] (overview) {G02 Traffic\\Overview};
\node[screen,below left=of overview] (perception) {G03 Perception and\\Intersection Detail};
\node[screen,below right=of overview] (emergency) {G04 Emergency\\Dispatch and Route};
\node[screen,below=23mm of overview] (analytics) {G05 Comparative\\Analytics};
\draw[arr] (config)--node[right]{validate and start}(overview);
\draw[arr,<->] (overview)--(perception);
\draw[arr,<->] (overview)--(emergency);
\draw[arr] (overview)--node[right]{completed run}(analytics);
\draw[arr] (analytics.west) -| (config.west);
\end{tikzpicture}
\caption{Proposed primary navigation flow between Manzil screens}
\label{fig:gui-navigation}
\end{figure}

The overview and perception wireframes provide complementary levels of detail. Figure \ref{fig:wire-overview} gives the operator a network-level view in which the selected run, controller, source mode, ambulance route, corridor lifecycle, and warnings remain visible together. This arrangement prevents a plausible map animation from hiding stale data or a fallback controller. Figure \ref{fig:wire-perception} narrows the view to one intersection and pairs a processed-video preview with model provenance, freshness, regional measures, confidence, and invalid-state indicators. The calibration field is explicit because pixel-derived queue proxies must not be mislabeled as physical queue length. Both designs use color only as a supporting cue; labels retain the meaning required for accessibility and printed evidence. The reader can therefore distinguish authoritative values, estimates, and warnings without relying on animation.

\begin{figure}[htbp]
\centering
\begin{tikzpicture}[x=1cm,y=1cm,font=\scriptsize]
\draw[rounded corners] (0,0) rectangle (14,7.2);
\fill[manzilblue] (0,6.5) rectangle (14,7.2); \node[text=white,anchor=west] at (.3,6.85) {MANZIL / Traffic Overview};
\draw[fill=manzilgray] (.3,.4) rectangle (9.3,6.2); \node at (4.8,5.9) {Network map: congestion, signals, ambulance and route};
\foreach \x/\y/\c in {1.4/4.8/manzilgreen,3.4/3.7/manzilamber,5.5/4.5/manzilred,7.6/2.6/manzilgreen}{\fill[\c] (\x,\y) circle (.16);}
\draw[very thick,manzilblue] (1.4,4.8)--(3.4,3.7)--(5.5,4.5)--(7.6,2.6);
\draw (9.6,4.7) rectangle (13.7,6.2); \node[align=left] at (11.65,5.45) {Run RUN-021\\Controller: Rule adaptive\\Source: SUMO / current};
\draw (9.6,2.7) rectangle (13.7,4.5); \node[align=left] at (11.65,3.6) {Emergency EV-01\\ETA: computed\\Corridor: TRANSITION\\Next: I-07};
\draw (9.6,.4) rectangle (13.7,2.5); \node[align=left] at (11.65,1.45) {Warnings / event timeline\\stale and fallback states\\remain visible};
\end{tikzpicture}
\caption{G02 traffic-overview screen wireframe}
\label{fig:wire-overview}
\end{figure}

\begin{figure}[htbp]
\centering
\begin{tikzpicture}[x=1cm,y=1cm,font=\scriptsize]
\draw[rounded corners] (0,0) rectangle (14,6.7);
\fill[manzilblue] (0,6) rectangle (14,6.7); \node[text=white,anchor=west] at (.3,6.35) {Intersection I-07 / Perception Detail};
\draw[fill=black!10] (.3,2.1) rectangle (8.4,5.7); \node at (4.35,5.35) {Processed video preview (when available)};
\draw[manzilgreen,thick] (1.2,2.8) rectangle (2.3,4.0); \node at (1.75,2.55) {track 17};
\draw[manzilamber,thick] (3.1,3.0) rectangle (4.4,4.2); \node at (3.75,2.75) {track 23};
\draw (8.7,4.1) rectangle (13.7,5.7); \node[align=left] at (11.2,4.9) {Source: selected clip\\Model: version / checksum\\Freshness: current or stale\\Calibration: unavailable};
\draw (8.7,2.1) rectangle (13.7,3.8); \node[align=left] at (11.2,2.95) {North: count --; queue proxy --\\East: count --; queue proxy --\\Illustrative fields, no live data};
\draw (.3,.3) rectangle (13.7,1.8); \node at (7,1.05) {Time-series area: count, occupancy, queue proxy, processing throughput};
\end{tikzpicture}
\caption{G03 perception and intersection-detail screen wireframe}
\label{fig:wire-perception}
\end{figure}

\begin{figure}[htbp]
\centering
\begin{tikzpicture}[x=1cm,y=1cm,font=\scriptsize]
\draw[rounded corners] (0,0) rectangle (14,6.3);
\fill[manzilblue] (0,5.6) rectangle (14,6.3); \node[text=white,anchor=west] at (.3,5.95) {Emergency Dispatch and Route};
\draw (.3,3.9) rectangle (5.0,5.3); \node[align=left] at (2.65,4.6) {Origin [select edge]\\Destination [select edge]\\Dispatch [button]};
\draw[fill=manzilgray] (5.3,.4) rectangle (13.7,5.3); \node at (9.5,5.0) {Map with route alternatives and current mission};
\draw[very thick,manzilgreen] (6.2,1.4)--(7.6,2.5)--(9.3,2.0)--(11.2,3.4)--(12.7,4.1);
\draw[thick,dashed,manzilamber] (6.2,1.4)--(8.5,3.6)--(10.4,4.2)--(12.7,4.1);
\draw (.3,1.9) rectangle (5,3.6); \node[align=left] at (2.65,2.75) {Selected cost: seconds\\Route version: 3\\Reroute threshold: config\\Reason: sustained improvement};
\draw (.3,.4) rectangle (5,1.6); \node[align=left] at (2.65,1.0) {Corridor lifecycle\\I-07 ACTIVE / I-08 REQUESTED};
\end{tikzpicture}
\caption{G04 emergency-dispatch and route screen wireframe}
\label{fig:wire-dispatch}
\end{figure}

Emergency operation and experiment setup require different decisions, so they are separated in Figures \ref{fig:wire-dispatch} and \ref{fig:wire-config}. The dispatch view shows validated endpoints, the current route version, its time-based cost, the reason for a revision, and the state of intersections in the bounded look-ahead corridor. It also permits a candidate route to be shown without implying that every alternative has been activated. The configuration view collects the network, demand, controller, seed, duration, incident, mission, and metric schema before a run begins. Its validation panel exposes unresolved identifiers, illegal phase plans, absent calibration, or checkpoint mismatch, and the start action remains unavailable until mandatory errors are resolved. The two wireframes therefore support auditability at different stages: one explains a live mission decision, while the other preserves the conditions under which that decision can later be compared. Neither view grants direct signal control, bypasses backend validation, or changes an active run's immutable configuration.

\begin{figure}[htbp]
\centering
\begin{tikzpicture}[x=1cm,y=1cm,font=\scriptsize]
\draw[rounded corners] (0,0) rectangle (14,6);
\fill[manzilblue] (0,5.3) rectangle (14,6); \node[text=white,anchor=west] at (.3,5.65) {Experiment Configuration};
\draw (.5,4.35) rectangle (8.4,4.85); \node[anchor=west] at (.7,4.6) {Network and demand profile};
\draw (.5,3.70) rectangle (8.4,4.20); \node[anchor=west] at (.7,3.95) {Controller and policy version};
\draw (.5,3.05) rectangle (8.4,3.55); \node[anchor=west] at (.7,3.30) {Seed, warm-up and duration};
\draw (.5,2.40) rectangle (8.4,2.90); \node[anchor=west] at (.7,2.65) {Incident and emergency mission};
\draw (.5,1.75) rectangle (8.4,2.25); \node[anchor=west] at (.7,2.00) {Metric schema and sample interval};
\draw (8.8,2.0) rectangle (13.5,4.85); \node[align=left] at (11.15,3.4) {Validation summary\\\checkmark IDs resolve\\\checkmark phase plan legal\\! video calibration absent\\! RL checkpoint mismatch};
\draw[fill=manzilgreen!25] (9.0,.7) rectangle (11.0,1.4); \node at (10,1.05) {Save draft};
\draw[fill=manzilblue!20] (11.3,.7) rectangle (13.3,1.4); \node at (12.3,1.05) {Start run};
\end{tikzpicture}
\caption{G01 experiment-configuration screen wireframe}
\label{fig:wire-config}
\end{figure}

\begin{figure}[htbp]
\centering
\begin{tikzpicture}[x=1cm,y=1cm,font=\scriptsize]
\draw[rounded corners] (0,0) rectangle (14,6);
\fill[manzilblue] (0,5.3) rectangle (14,6); \node[text=white,anchor=west] at (.3,5.65) {Comparative Analytics};
\draw (.4,4.15) rectangle (13.6,5.0); \node at (7,4.58) {Run selector / comparability warning / metric definition};
\draw (.4,.4) rectangle (8.7,3.85); \node at (4.55,3.55) {Distribution with mean and variability};
\draw[fill=manzilblue!20] (1.3,.8) rectangle (2.4,2.2); \draw[fill=manzilgreen!20] (3.2,.8) rectangle (4.3,2.2); \draw[fill=manzilamber!30] (5.1,.8) rectangle (6.2,2.2);
\node at (1.85,.6) {Fixed}; \node at (3.75,.6) {Rule}; \node at (5.65,.6) {RL}; \node at (4.0,3.15) {Illustrative layout; no measured values};
\draw (9.0,.4) rectangle (13.6,3.85); \node[align=left] at (11.3,2.1) {Configuration differences\\Seed coverage\\Failed runs\\Export data + manifest};
\end{tikzpicture}
\caption{G05 comparative-analytics screen wireframe}
\label{fig:wire-analytics}
\end{figure}

Figure \ref{fig:wire-analytics} represents the final evidence-review workflow. Run selection is paired with a comparability warning and metric definition so that two visually adjacent bars are not assumed to represent matched experiments. The chart area is schematic and contains no claimed result; an implemented view would present the chosen distribution, central tendency, variability, and sample count from retained run records. The adjoining panel exposes configuration differences, seed coverage, failed runs, and export controls. This design responds directly to the methodology requirement that failed or incompatible runs must not disappear from a favorable summary. It also keeps measurement units and metric versions close to the visualization, reducing the risk that a screenshot is detached from its experimental meaning. Together, the five wireframes define the minimum information architecture for the prototype while leaving visual polish and specific chart libraries open to implementation. Exported graphics must retain the same metric and run identifiers.

\section{Database Design}
\label{sec:data-design}

\subsection{Entity Relationship Diagram}

Persistent data is reserved for configuration, observations, decisions, missions, and outcomes that require audit or comparison. The live road graph, TraCI subscription, and frame buffers remain runtime state, while raw video is retained as a checksummed external artifact. Following the notation used in the reference report, Figure \ref{fig:er} presents the complete conceptual model with blue entity rectangles, green attribute ovals, amber relationship diamonds, and cardinalities beside relationship lines. The dashed regions separate experiment configuration and evidence, observation and control records, and emergency routing and priority. Figures \ref{fig:er-experiment}, \ref{fig:er-observation}, and \ref{fig:er-emergency} enlarge those regions so keys, attributes, and named relationships remain readable at normal print size. Primary keys are underlined, foreign keys are marked FK, and optional ownership is expressed explicitly. The diagrams describe a proposed logical schema and do not claim that a production database has already been populated.

\begin{figure}[htbp]
\centering
\begin{tikzpicture}[x=1cm,y=1cm,font=\tiny,
entity/.style={draw=manzilblue,fill=manzilblue!72,text=white,rectangle,minimum width=1.65cm,minimum height=.48cm,align=center},
attribute/.style={draw=manzilgreen,fill=manzilgreen!13,ellipse,minimum width=1.55cm,minimum height=.36cm,align=center},
relationship/.style={draw=manzilamber,fill=manzilamber!35,diamond,aspect=2.2,inner sep=1pt,align=center},
link/.style={draw=black!70,thin}]
\draw[dashed,manzilgreen,rounded corners] (-.2,.1) rectangle (6.8,4.8);
\node[anchor=north west,font=\scriptsize] at (0,4.65) {Experiment Configuration and Evidence};
\node[entity] (sc) at (1.0,3.25) {Scenario};
\node[attribute] (scid) at (1.0,4.05) {\underline{scenario\_id}};
\node[entity] (run) at (4.0,3.25) {ExperimentRun};
\node[attribute] (runid) at (3.4,4.05) {\underline{run\_id}};
\node[attribute] (seed) at (5.0,4.05) {seed};
\node[relationship] (defines) at (2.5,3.25) {defines};
\node[entity] (metric) at (1.0,1.35) {MetricSample};
\node[attribute] (metricid) at (1.0,.62) {\underline{metric\_id}};
\node[entity] (artifact) at (4.0,1.35) {Artifact};
\node[attribute] (artifactid) at (4.0,.62) {\underline{artifact\_id}};
\node[relationship] (records) at (2.5,2.25) {records};
\draw[link] (sc)--(scid); \draw[link] (run)--(runid); \draw[link] (run)--(seed);
\draw[link] (sc)--node[above]{1}(defines)--node[above]{0..*}(run);
\draw[link] (run)--(records); \draw[link] (records)-|node[pos=.72,above]{0..*}(metric);
\draw[link] (records)-|node[pos=.72,above]{0..*}(artifact);
\draw[link] (metric)--(metricid); \draw[link] (artifact)--(artifactid);

\draw[dashed,manzilgreen,rounded corners] (7.1,.1) rectangle (14.1,4.8);
\node[anchor=north west,font=\scriptsize] at (7.3,4.65) {Observation and Control Evidence};
\node[entity] (video) at (8.2,3.25) {VideoJob};
\node[attribute] (videoid) at (8.2,4.05) {\underline{video\_job\_id}};
\node[entity] (obs) at (11.0,3.25) {Observation};
\node[attribute] (obsid) at (11.0,4.05) {\underline{observation\_id}};
\node[entity] (decision) at (11.0,1.35) {ControlDecision};
\node[attribute] (decisionid) at (11.0,.62) {\underline{decision\_id}};
\node[entity] (runobs) at (8.2,1.35) {ExperimentRun};
\node[relationship] (ownsobs) at (9.6,1.35) {owns};
\node[relationship] (processes) at (9.6,3.25) {produces};
\node[relationship] (informs) at (11.0,2.25) {informs};
\draw[link] (video)--(videoid); \draw[link] (obs)--(obsid); \draw[link] (decision)--(decisionid);
\draw[link] (video)--node[above]{1}(processes)--node[above]{0..*}(obs);
\draw[link] (obs)--node[right]{1}(informs)--node[right]{0..*}(decision);
\draw[link] (runobs)--node[above]{1}(ownsobs)--node[above]{0..*}(decision);

\draw[dashed,manzilgreen,rounded corners] (-.2,-5.0) rectangle (14.1,-.25);
\node[anchor=north west,font=\scriptsize] at (0,-.4) {Emergency Routing and Priority};
\node[entity] (runem) at (.8,-2.2) {ExperimentRun};
\node[relationship] (ownsem) at (2.0,-2.2) {owns};
\node[entity] (mission) at (3.3,-2.2) {EmergencyMission};
\node[attribute] (missionid) at (3.3,-1.35) {\underline{mission\_id}};
\node[relationship] (dispatches) at (5.25,-2.2) {versions};
\node[entity] (route) at (7.2,-2.2) {RouteVersion};
\node[attribute] (routeid) at (6.65,-1.35) {\underline{route\_id}};
\node[attribute] (version) at (8.2,-1.35) {version};
\node[relationship] (authorises) at (9.3,-2.2) {authorises};
\node[entity] (priority) at (11.5,-2.2) {PriorityEvent};
\node[attribute] (eventid) at (11.5,-1.35) {\underline{event\_id}};
\draw[link] (mission)--(missionid); \draw[link] (route)--(routeid); \draw[link] (route)--(version); \draw[link] (priority)--(eventid);
\draw[link] (runem)--node[above]{1}(ownsem)--node[above]{0..*}(mission);
\draw[link] (mission)--node[above]{1}(dispatches)--node[above]{1..*}(route);
\draw[link] (route)--node[above]{1}(authorises)--node[above]{0..*}(priority);
\end{tikzpicture}
\caption{Complete grouped conceptual ER model using entity, attribute, and relationship symbols}
\label{fig:er}
\end{figure}

Figure \ref{fig:er-experiment} enlarges the experiment-configuration region of the overview. A Scenario defines immutable network and demand versions and may be reused by several ExperimentRun records. Each run fixes its seed, controller version, status, and configuration at start, then owns its MetricSample and Artifact records. The records relationship is deliberately named because metrics and exported files are evidence produced under one run rather than free-standing data. Metric samples retain a definition version, interval, value, and unit, while artifacts retain a type, URI, checksum, and provenance. These associations prevent results from different seeds or controller configurations from being combined without their originating run. They also permit an evaluator to trace a displayed chart or exported file back to the configuration that produced it.

\begin{figure}[htbp]
\centering
\begin{tikzpicture}[x=1cm,y=1cm,font=\scriptsize,
entity/.style={draw=manzilblue,fill=manzilblue!72,text=white,rectangle,minimum width=2.2cm,minimum height=.62cm,align=center},
attribute/.style={draw=manzilgreen,fill=manzilgreen!13,ellipse,minimum width=2.15cm,minimum height=.45cm,align=center},
relationship/.style={draw=manzilamber,fill=manzilamber!35,diamond,aspect=2.2,inner sep=1.5pt,align=center},
link/.style={draw=black!70,thin}]
\draw[manzilgreen,rounded corners] (-.2,-5.1) rectangle (14.1,3.3);
\node[anchor=north west,font=\small] at (0,3.1) {Experiment Configuration and Evidence};
\node[entity] (sc) at (1.6,1.2) {Scenario};
\node[attribute] (scid) at (.5,2.25) {\underline{scenario\_id}};
\node[attribute] (network) at (2.75,2.25) {network\_version};
\node[attribute] (demand) at (1.6,.15) {demand\_version};
\node[relationship] (defines) at (4.0,1.2) {defines};
\node[entity] (run) at (6.3,1.2) {ExperimentRun};
\node[attribute] (runid) at (5.0,2.25) {\underline{run\_id}};
\node[attribute] (scenariofk) at (7.55,2.25) {scenario\_id (FK)};
\node[attribute] (seed) at (5.0,.15) {seed};
\node[attribute] (controller) at (7.55,.15) {controller\_version};
\node[attribute] (status) at (6.3,-.75) {status};
\draw[link] (sc)--(scid); \draw[link] (sc)--(network); \draw[link] (sc)--(demand);
\draw[link] (run)--(runid); \draw[link] (run)--(scenariofk); \draw[link] (run)--(seed); \draw[link] (run)--(controller); \draw[link] (run)--(status);
\draw[link] (sc)--node[above]{1}(defines)--node[above]{0..*}(run);
\node[relationship] (records) at (8.8,-1.7) {records};
\draw[link] (run)--node[right]{1}(records);
\node[entity] (metric) at (6.2,-3.3) {MetricSample};
\node[attribute] (metricid) at (4.8,-4.35) {\underline{metric\_id}};
\node[attribute] (metricrun) at (7.5,-4.35) {run\_id (FK)};
\node[attribute] (metricname) at (4.6,-2.35) {metric\_name};
\node[attribute] (metricvalue) at (7.8,-2.35) {value, unit};
\node[entity] (artifact) at (11.5,-3.3) {Artifact};
\node[attribute] (artifactid) at (10.2,-4.35) {\underline{artifact\_id}};
\node[attribute] (artifactrun) at (12.8,-4.35) {run/job\_id (FK)};
\node[attribute] (uri) at (10.2,-2.35) {URI};
\node[attribute] (checksum) at (12.8,-2.35) {checksum};
\draw[link] (records)-|node[pos=.72,above]{0..*}(metric); \draw[link] (records)-|node[pos=.72,above]{0..*}(artifact);
\draw[link] (metric)--(metricid); \draw[link] (metric)--(metricrun); \draw[link] (metric)--(metricname); \draw[link] (metric)--(metricvalue);
\draw[link] (artifact)--(artifactid); \draw[link] (artifact)--(artifactrun); \draw[link] (artifact)--(uri); \draw[link] (artifact)--(checksum);
\end{tikzpicture}
\caption{Experiment scenarios, runs, metrics, and artifacts with conceptual attributes}
\label{fig:er-experiment}
\end{figure}

Figure \ref{fig:er-observation} separates the two permitted observation owners and connects observations to control evidence. Simulation observations reference an ExperimentRun; offline video observations reference a VideoJob. Exactly one owner is present for each observation, so a record cannot be presented simultaneously as simulator truth and video-derived evidence. The Observation entity carries source mode, entity identifier, effective and validity times, values, units, and confidence. A ControlDecision refers to the run and to the observation context used by the controller, while preserving requested and applied phases and any safety reason. The named informs relationship expresses causal use rather than simple co-location. This structure supports stale-data checks and permits the evaluator to distinguish a controller error from an unavailable observation, rejected action, or safety transformation.

\begin{figure}[htbp]
\centering
\begin{tikzpicture}[x=1cm,y=1cm,font=\scriptsize,
entity/.style={draw=manzilblue,fill=manzilblue!72,text=white,rectangle,minimum width=2.25cm,minimum height=.62cm,align=center},
attribute/.style={draw=manzilgreen,fill=manzilgreen!13,ellipse,minimum width=2.15cm,minimum height=.45cm,align=center},
relationship/.style={draw=manzilamber,fill=manzilamber!35,diamond,aspect=2.2,inner sep=1.5pt,align=center},
link/.style={draw=black!70,thin}]
\draw[manzilgreen,rounded corners] (-.2,-5.0) rectangle (14.1,3.25);
\node[anchor=north west,font=\small] at (0,3.05) {Observation and Control Evidence};
\node[entity] (video) at (1.4,1.15) {VideoJob};
\node[attribute] (videoid) at (.45,2.15) {\underline{video\_job\_id}};
\node[attribute] (clip) at (2.65,2.15) {clip\_checksum};
\node[attribute] (model) at (1.4,.15) {model\_version};
\node[relationship] (produces) at (4.0,1.15) {produces};
\node[entity] (obs) at (6.5,1.15) {Observation};
\node[attribute] (obsid) at (5.25,2.15) {\underline{observation\_id}};
\node[attribute] (owner) at (7.85,2.15) {run/job\_id (FK)};
\node[attribute] (source) at (5.25,.15) {source\_mode};
\node[attribute] (valid) at (7.85,.15) {valid\_until};
\node[attribute] (values) at (4.9,-.7) {values, units, confidence};
\node[relationship] (informs) at (6.5,-2.0) {informs};
\node[entity] (decision) at (10.7,-2.0) {ControlDecision};
\node[attribute] (decisionid) at (9.4,-3.05) {\underline{decision\_id}};
\node[attribute] (decisionrun) at (12.2,-3.05) {run\_id (FK)};
\node[attribute] (requested) at (9.3,-1.0) {requested\_phase};
\node[attribute] (applied) at (12.2,-1.0) {applied\_phase};
\node[attribute] (reason) at (10.7,-4.2) {reason, effective\_step};
\node[entity] (run) at (1.4,-3.55) {ExperimentRun};
\node[relationship] (owns) at (4.0,-3.55) {owns};
\draw[link] (video)--(videoid); \draw[link] (video)--(clip); \draw[link] (video)--(model);
\draw[link] (video)--node[above]{1}(produces)--node[above]{0..*}(obs);
\draw[link] (run)--node[above]{1}(owns)--node[above]{0..*}(decision);
\draw[link] (obs) -- node[right]{1}(informs) -- node[above]{0..*}(decision);
\draw[link] (obs)--(obsid); \draw[link] (obs)--(owner); \draw[link] (obs)--(source); \draw[link] (obs)--(valid); \draw[link] (obs)--(values);
\draw[link] (decision)--(decisionid); \draw[link] (decision)--(decisionrun); \draw[link] (decision)--(requested); \draw[link] (decision)--(applied); \draw[link] (decision)--(reason);
\end{tikzpicture}
\caption{Observation ownership and control-decision evidence with conceptual attributes}
\label{fig:er-observation}
\end{figure}

Figure \ref{fig:er-emergency} enlarges the emergency-control chain. A run can own several mission attempts over its lifetime, although the core scope permits only one active mission at a time. Each accepted dispatch creates an EmergencyMission with vehicle, origin, destination, and status. An accepted route or reroute creates a new immutable RouteVersion rather than overwriting the previous path. The versions relationship is one-to-many and the version number increases within its mission. PriorityEvent records then reference the exact route version that authorised a request, safe transition, active-priority interval, passage, cancellation, or recovery state. This structure allows obsolete priority work to be cancelled after rerouting and later explains which path justified every intersection action. The ER notation therefore reflects the same causal ownership enforced by the corridor state machine.

\begin{figure}[htbp]
\centering
\begin{tikzpicture}[x=1cm,y=1cm,font=\scriptsize,
entity/.style={draw=manzilblue,fill=manzilblue!72,text=white,rectangle,minimum width=2.25cm,minimum height=.62cm,align=center},
attribute/.style={draw=manzilgreen,fill=manzilgreen!13,ellipse,minimum width=2.15cm,minimum height=.45cm,align=center},
relationship/.style={draw=manzilamber,fill=manzilamber!35,diamond,aspect=2.2,inner sep=1.5pt,align=center},
link/.style={draw=black!70,thin}]
\draw[manzilgreen,rounded corners] (-.2,-5.0) rectangle (14.1,3.25);
\node[anchor=north west,font=\small] at (0,3.05) {Emergency Routing and Priority};
\node[entity] (run) at (1.25,1.25) {ExperimentRun};
\node[relationship] (owns) at (3.2,1.25) {owns};
\node[entity] (mission) at (5.3,1.25) {EmergencyMission};
\node[attribute] (missionid) at (4.25,2.25) {\underline{mission\_id}};
\node[attribute] (missionrun) at (6.55,2.25) {run\_id (FK)};
\node[attribute] (vehicle) at (7.85,1.45) {vehicle\_id};
\node[attribute] (endpoints) at (7.85,.55) {origin, destination};
\node[relationship] (versions) at (5.3,-.05) {versions};
\node[entity] (route) at (5.3,-1.3) {RouteVersion};
\node[attribute] (routeid) at (2.8,-.85) {\underline{route\_id}};
\node[attribute] (missionfk) at (2.8,-1.75) {mission\_id (FK)};
\node[attribute] (version) at (7.85,-.85) {version};
\node[attribute] (edges) at (8.3,-1.75) {edges, cost, trigger};
\node[relationship] (authorises) at (5.3,-2.55) {authorises};
\node[entity] (event) at (5.3,-3.75) {PriorityEvent};
\node[attribute] (eventid) at (2.8,-3.35) {\underline{event\_id}};
\node[attribute] (routefk) at (2.8,-4.3) {route\_id (FK)};
\node[attribute] (intersection) at (8.0,-3.35) {intersection\_id};
\node[attribute] (state) at (8.0,-4.3) {state, effective\_step};
\draw[link] (run)--node[above]{1}(owns)--node[above]{0..*}(mission);
\draw[link] (mission)--(missionid); \draw[link] (mission)--(missionrun); \draw[link] (mission)--(vehicle); \draw[link] (mission)--(endpoints);
\draw[link] (mission)--node[right]{1}(versions)--node[right]{1..*}(route);
\draw[link] (route)--(routeid); \draw[link] (route)--(missionfk); \draw[link] (route)--(version); \draw[link] (route)--(edges);
\draw[link] (route)--node[right]{1}(authorises)--node[right]{0..*}(event);
\draw[link] (event)--(eventid); \draw[link] (event)--(routefk); \draw[link] (event)--(intersection); \draw[link] (event)--(state);
\end{tikzpicture}
\caption{Emergency missions, route versions, and priority events with conceptual attributes}
\label{fig:er-emergency}
\end{figure}

\subsection{Data Dictionary}

Table \ref{tab:data-dictionary} defines the principal identifiers and integrity constraints represented in the ER diagrams. It adds the provenance and validity rules that cannot be communicated by cardinality alone, including immutable run configuration, mandatory observation source and units, monotonically increasing route versions, and metric-version compatibility. The dictionary also distinguishes large external artifacts from relational records so videos, models, and logs can remain checksummed files without becoming database blobs. Together, the diagram and dictionary establish the proposed storage contract for implementation and evaluation.

\begin{longtable}{P{0.16\textwidth}P{0.28\textwidth}P{0.45\textwidth}}
\caption{Data dictionary}\label{tab:data-dictionary}\\
\toprule
\textbf{Entity} & \textbf{Key fields} & \textbf{Integrity rule}\\
\midrule\endfirsthead
\toprule
\textbf{Entity} & \textbf{Key fields} & \textbf{Integrity rule}\\
\midrule\endhead
Scenario & scenario\_id, network\_version, demand\_version & References immutable network and demand manifests.\\
Run & run\_id, scenario\_id, seed, controller\_version, status & Configuration becomes immutable after start; terminal reason required.\\
VideoJob & job\_id, clip\_checksum, model\_version, status & Clip provenance and model version become immutable after processing begins.\\
Observation & run\_id or job\_id, source\_mode, entity\_id, observed\_at, valid\_until, values, units & Exactly one owner is present; source and units are mandatory and validity cannot be inferred.\\
ControlDecision & run\_id, intersection\_id, requested\_phase, applied\_phase, effective\_step, reason & Requested and applied actions are separate fields.\\
Mission & mission\_id, run\_id, vehicle\_id, origin, destination, status & Only one active core mission per run.\\
RouteVersion & mission\_id, version, edge\_list, cost\_seconds, trigger, created\_at & Version increases monotonically; path begins at valid current position.\\
PriorityEvent & mission\_id, route\_version, intersection\_id, state, requested/applied time & References the route version that justified priority.\\
MetricSample & run\_id, metric\_name, metric\_version, value, unit, interval & Comparisons require matching metric name, version, and unit.\\
Artifact & run/job\_id, type, URI, checksum, license/provenance & Content is addressable and provenance is retained.\\
\bottomrule
\end{longtable}

\section{Risk Analysis}

Table \ref{tab:risks} records technical and research risks that could invalidate the prototype or its conclusions. Likelihood and impact are qualitative planning judgments rather than measured probabilities. The highest-consequence risks concern invalid signal phases, irreproducible results, unsuitable map imports, and evidence that cannot be legally or reliably used. Each mitigation is paired with a bounded fallback so that the team can preserve an honest deliverable if a preferred component fails. For example, a documented synthetic network is preferable to an unverified real-map extract, and a negative reinforcement-learning result is preferable to tuning on the final test scenarios. The register also treats timing drift, route oscillation, occlusion, and reward imbalance as system-level concerns rather than isolated coding defects. These risks are revisited through invariant tests, scenario checks, provenance records, and the controlled comparisons defined in Chapter \ref{ch:method}; unresolved items must remain visible in the final limitations rather than being converted into unsupported performance claims.

\begin{longtable}{P{0.20\textwidth}P{0.16\textwidth}P{0.34\textwidth}P{0.19\textwidth}}
\caption{Risk register}\label{tab:risks}\\
\toprule
\textbf{Risk} & \textbf{Rating} & \textbf{Mitigation} & \textbf{Fallback}\\
\midrule\endfirsthead
Map import lacks valid signals/turns & Medium likelihood; high impact & Inspect and clean a small extract early; validate routes and phases & Use a documented synthetic grid\\
Local video lacks labels or permission & Medium likelihood; high impact & Use licensed public sequences and record provenance & Evaluate public data only; state domain gap\\
Occlusion causes unstable counts & High likelihood; medium impact & Use region rules, temporal aggregation, and tracker comparison & Limit claims to labelled conditions\\
RL fails to converge or generalize & Medium likelihood; high impact & Use stable baselines, bounded actions, and training-only tuning & Report a negative or partial finding\\
Reward improves one metric unfairly & Medium likelihood; high impact & Evaluate multiple metrics and enforce starvation constraints & Reject a policy that violates fairness checks\\
Rerouting oscillates & Medium likelihood; medium impact & Apply a benefit threshold, persistence, cooldown, and closure exception & Retain the current feasible route\\
Preemption leaves unsafe phase & Low likelihood; very high impact & Use one state machine and test transition invariants & Cancel and enter a safe recovery path\\
Service timing drifts from simulation & Medium likelihood; high impact & Use one step owner, simulation timestamps, and bounded queues & Run synchronously without live perception\\
Insufficient compute & Medium likelihood; medium impact & Profile early; reduce network/model size and train offline & Demonstrate recorded inference and compact scenarios\\
Results are not reproducible & Medium likelihood; high impact & Retain immutable manifests, seeds, versioned metrics, and exports & Exclude untraceable runs\\
\bottomrule
\end{longtable}

\section{Conclusion}


The specification centers on traceability and failure visibility rather than on a collection of loosely connected features. It requires one simulation-clock owner, legal signal transitions, stable route revision, explicit observation freshness, immutable run configuration, and comparisons that cannot hide incompatible settings or failed trials. Functional requirements define what the operator and services must observe, while non-functional requirements constrain timing, ownership, privacy, auditability, and recovery. The use cases exercise both successful and failed paths, the wireframes show how critical state and provenance remain visible, and the data model preserves the route and decision history needed for later analysis. The risk controls connect those requirements to bounded fallbacks and planned verification. Chapter \ref{ch:method} converts the specification into a staged research protocol, and Chapter \ref{ch:design} assigns each mutable resource and interface to a clear architectural owner. Implementation status will continue to be reported separately so that a specified feature is never presented as completed evidence.
